\documentclass[11pt]{amsart}
\usepackage[margin=1.1in]{geometry}
\usepackage{amsmath,amssymb,amsthm}
\usepackage{booktabs}
\usepackage[colorlinks=true,linkcolor=blue,citecolor=blue,urlcolor=blue]{hyperref}

\newtheorem{theorem}{Theorem}[section]
\newtheorem{lemma}[theorem]{Lemma}
\newtheorem{proposition}[theorem]{Proposition}
\newtheorem{corollary}[theorem]{Corollary}
\theoremstyle{definition}
\newtheorem{definition}[theorem]{Definition}
\newtheorem{example}[theorem]{Example}
\newtheorem{axiom}{Axiom}
\theoremstyle{remark}
\newtheorem{remark}[theorem]{Remark}

\newcommand{\Ran}{\operatorname{Ran}}
\newcommand{\sgn}{\operatorname{sgn}}
\newcommand{\artanh}{\operatorname{artanh}}
\newcommand{\csch}{\operatorname{csch}}
\newcommand{\Hil}{\mathcal H}
\newcommand{\eps}{\varepsilon}
\newcommand{\one}{\mathbf 1}
\emergencystretch=3em

\begin{document}

\title[Schr\"odinger as the $\alpha\to1^-$ collapse of the CAL]{Schr\"odinger evolution as the $\alpha\to1^-$ collapse of the Causal Accumulation Law: derived unitarity at the free Galilean edge}

\author{Sam Leizerman}
\address{LIOR Research Initiative}
\email{samleizerman@outlook.com}
\thanks{ORCID: 0009-0000-0133-2291. Companion to \emph{The Causal Accumulation Law: The Four Corners of $\alpha$, Algebraic Closure, and the Standard Model}~\cite{CAL} and to \emph{The $\alpha=1$ Unitary Skeleton and the Emergence of Schr\"odinger Evolution in the $\delta$-Kernel Limit of the CAL}~\cite{Skeleton}, whose unitarity posit it replaces. Dated November 11, 2026.}

\begin{abstract}
We impose the Causal Accumulation Law (CAL) as a field equation, $T=0$, with a power-law memory under causal Weyl integration and the channel orientation of the parent framework carried by the memory channel. As the fractional order approaches $\alpha=1$ from below, the field equation collapses to the Schr\"odinger equation.

Unitarity is derived, not posited. For every $\alpha<1$, the bounded solutions on the whole time line form a strongly continuous unitary group, because the channel phase cancels the phase of the Weyl symbol on one sign of frequency. This replaces the unitarity posit of the companion paper.

The limit is one-sided for three reasons. The Weyl integral converges on oscillating carriers only for $\alpha<1$. The memory-generated coefficients change sign across $\alpha=1$. And from above, the response to a source either runs away at rate $\pi\lvert1-\alpha\rvert\,\lvert E\rvert/\hbar$ or precedes its source, which is the Abraham--Lorentz--Dirac dichotomy.

The kernel-blind skeleton at $\alpha=1$ and the Schr\"odinger edge are two scalings of one limit. With the kernel's time unit held fixed, the memory rate vanishes and the past disconnects even from below. The edge exists only if the memory rate $(1-\alpha)/\tau_c$ is held fixed. In that limit the kernel becomes a flat accumulator, not a $\delta$.

Neither of the framework's two Markovian regimes yields unitary Schr\"odinger evolution by itself.
\begin{itemize}
\item At $\alpha=1$ the past is gone, so there is no generator.
\item A normalizable kernel, including every $\delta$-like family, has no stationary band at all. The plane-wave equation needs $\lvert\hat w\rvert=\alpha/(\eps\lvert\lambda\rvert)$, while $\lvert\hat w\rvert\le1$.
\end{itemize}
The power law succeeds because it is not normalizable, and because the phase of its symbol is the channel angle. The edge therefore has a memory law with a Markovian flow.

The edge is Galilean and non-interacting.
\begin{itemize}
\item The Lorentz signature belongs to the $\alpha=2$ closure, so the edge has no light cone. This is consistent with the instantaneous spatial reach of the Schr\"odinger propagator.
\item Its generator is built from one-way transports. It therefore gives the free Schr\"odinger form, at most in a fixed one-way background.
\item Interaction and entanglement need the $\alpha=2$ closure of a composite carrier.
\end{itemize}

Numerical runs of the field equation, using Weyl integration only, confirm each statement.
\end{abstract}

\maketitle
\tableofcontents

%======================================================================
\section{Scope and claim}

The companion paper~\cite{Skeleton} obtains the Schr\"odinger form at the edge of the $(\alpha,\text{kernel})$ table adjoining $\alpha=1$ from below, but it consumes an input: Posit~1 assumes a unitary flow $U(t)=e^{i\theta(t)G}$ and then defines $\hat H\equiv i\hbar\dot U U^{-1}$. That step establishes the form of the equation only after unitarity and a generator have been assumed. The present paper removes the posit. The accumulation law itself, read as a field equation, collapses to the Schr\"odinger equation, and unitarity, the one-sidedness of the limit, and the mass formula follow.

\begin{quote}\itshape
Impose $T=0$ with a power-law memory under causal Weyl integration, and let the memory channel carry the orientation $e^{-i\sigma\alpha\pi/2}$ of Axiom~3. For every $\alpha<1$, the bounded solutions on the whole time line evolve by a unitary group (Theorem~\ref{thm:eternal}). As $\alpha\to1^-$ with the memory rate held fixed, they converge to Schr\"odinger evolution (Theorem~\ref{thm:collapse}). From above, the response to a source runs away or precedes the source (Theorem~\ref{thm:onesided}). Neither Markovian regime of the framework does this on its own: normalizable kernels, $\delta$-like ones included, have no stationary band (Theorem~\ref{thm:normalizable}). The edge is Galilean and carries no interaction (\S\ref{sec:galilean}--\ref{sec:nointeraction}).
\end{quote}

All dynamics in this paper use Weyl integration over the whole past, which is the integral the accumulation law writes. Initial-value (Caputo) problems are not used. A Caputo problem is the Weyl derivative of a history frozen at its initial value for all earlier times. Such a frozen past has zero-frequency content, on which the Weyl integral of the field equation diverges, so a Caputo problem is not a solution of Axiom~\ref{ax:FE}. Starting a solution is instead modelled by a source acting on a zero past (Proposition~\ref{prop:switch}).

\emph{Bookkeeping.} The inputs are:
\begin{itemize}
\item the field equation (Axiom~\ref{ax:FE}), which replaces Posit~1 of~\cite{Skeleton};
\item the kernel class (Definition~\ref{def:kernel}): a power law in the lag with a lag-independent spatial profile;
\item the framework's assignment of the Lorentz signature to the $\alpha=2$ closure~\cite{CAL}, which makes the causal structure at the edge time ordering only;
\item the double scaling (Definition~\ref{def:double}), in which the memory rate is held fixed.
\end{itemize}
Given these inputs, the following are derived:
\begin{itemize}
\item Lemmas~\ref{lem:weyl}--\ref{lem:ops};
\item Theorems~\ref{thm:eternal}, \ref{thm:collapse}, \ref{thm:onesided} and~\ref{thm:normalizable};
\item Propositions~\ref{prop:switch} and~\ref{prop:rate}, and Corollary~\ref{cor:scalings}.
\end{itemize}
The following are not claimed:
\begin{itemize}
\item a value for the mass scale;
\item interaction in the reciprocal sense: back-reaction, two-body coupling, entangling dynamics, or a gauge field strength generated by the carrier. These belong to the $\alpha=2$ closure (Remark~\ref{rem:nointeraction});
\item Lorentz structure, which also belongs to $\alpha=2$;
\item anything about measurement~\cite{CAL,Born}.
\end{itemize}

%======================================================================
\section{Definitions}

Throughout, $\eps\equiv1-\alpha$. Complex powers take the principal branch, and $E_c\equiv\hbar/\tau_c$.

\begin{definition}[Causal Weyl integral and admissibility]\label{def:weyl}
For a bounded $f:\mathbb R\to\mathcal V$ and $\alpha\in(0,1)$, the causal (left-sided) Weyl integral is
\begin{equation}\label{eq:weyl}
W^\alpha f(t)=\frac1{\Gamma(\alpha)}\int_0^\infty s^{\alpha-1}f(t-s)\,ds .
\end{equation}
It integrates over the past only. The right-sided integral $\int_t^\infty$, which some texts also call Weyl's~\cite{SKM}, integrates over the future and is excluded.

A function $f$ is \emph{admissible} if its time Fourier transform, as a tempered distribution, has no support at zero frequency. On admissible functions, $D^\alpha$ denotes the inverse of $W^\alpha$. On a carrier $e^{-iEt/\hbar}$ the symbols are $W^\alpha\mapsto(-iE/\hbar)^{-\alpha}$ and $D^\alpha\mapsto(-iE/\hbar)^{\alpha}$.
\end{definition}

\begin{definition}[Memory kernel]\label{def:kernel}
The memory kernel is
\[
k(\tau,y)=\tau_c^{-\alpha}\,\frac{\tau^{\alpha-1}}{\Gamma(\alpha)}\,g(y),
\]
dressed by the real transport $\mathcal P$. Here $\tau>0$ is the lag, and $g\ge0$ is a normalized, isotropic spatial profile that does not depend on the lag.

Below the $\alpha=2$ closure the causal structure is time ordering only (\S\ref{sec:galilean}). So the past $J^-(x)$ at the edge is $\{t'<t\}$, and $g$ is not required to shrink with the lag.

The \emph{memory rate} at lag $\tau$ is $\rho_\eps(\tau)\equiv(1-\alpha)\,\tau_c^{-\alpha}\tau^{\alpha-1}/\Gamma(\alpha)$.
\end{definition}

\begin{definition}[Channel orientation]\label{def:orient}
The orientation is $\sigma\in\{+1,-1\}$. By Axiom~3 of~\cite{Skeleton}, the memory channel carries the phase $e^{-i\sigma\theta}$ relative to the present channel, with $\theta=\alpha\pi/2$.
\end{definition}

\begin{definition}[Hilbert sector]\label{def:sector}
The spaces $\mathcal V$ and $\Hil$, the projection $\Pi_\Hil$ and the map $f_{\rm map}$ are those of~\cite[Defs.~2.3--2.5]{Skeleton}.

On the spinless, gauge-neutral sector $\Hil=L^2(\mathbb R^3)$, translation and rotation invariance give the projected memory the spatial form $\hat R_\Hil=r_0\nabla^2+r_1$~\cite[Thm.~6.1(ii)]{Skeleton}. Here $r_1$ is the coincidence weight of the fiber factors and $r_0=r_1\sigma_g^2/6$, with $\sigma_g^2$ the second moment of $g$. For $r_1>0$ this gives $r_0>0$.
\end{definition}

\begin{definition}[Double scaling]\label{def:double}
Take $\eps\to0^+$ with $\Lambda\equiv\eps\hat R_\Hil$ held fixed.

The limiting generator $H=\sigma E_c\Lambda=\sigma\hbar\,(\eps/\tau_c)\,\hat R_\Hil$ depends on $\eps$, $\tau_c$ and $\hat R_\Hil$ only through $(\eps/\tau_c)\hat R_\Hil$. So for bounded fiber factors this is the same limit as holding the memory rate $\eps/\tau_c$ fixed with $\hat R_\Hil$ fixed (Proposition~\ref{prop:rate}(2)). The two choices differ only in the approach, $O(\eps)$ against $O(\eps\ln\eps)$. If $\tau_c$ and $\hat R_\Hil$ are both held fixed instead, the past disconnects (Proposition~\ref{prop:rate}(1)).
\end{definition}

%======================================================================
\section{Axioms and the field equation}

\begin{axiom}[Causal accumulation law; C8 of~\cite{CAL}]\label{ax:CAL}
\[
T(x)=\alpha J(x)-(1-\alpha)\int_{J^-(x)}k(\tau;x,x')\,\mathcal P(x,x')\,J(x')\,d^4x'.
\]
The kernel $k$ and the transport $\mathcal P$ are real. $J$ takes values in the biquaternion carrier $\mathbb C\otimes\mathbb H$.
\end{axiom}

\begin{remark}
\cite{Skeleton} takes $J$ to be real. That is incompatible with its own phase channel $B\exp(iS/\hbar_{\rm eff})$, and with any solution proportional to $e^{-iEt/\hbar}$. The reality that~\cite[Lemma~4.3]{Skeleton} actually uses is the reality of $\hat R$, and that follows from the reality of $k$ and $\mathcal P$.
\end{remark}

\begin{axiom}[Two channels and channel geometry]
These are as in~\cite[Axioms~2--3]{Skeleton}: the channels carry relative orientation $e^{i\theta}$ with $\theta=\alpha\pi/2$.
\end{axiom}

\begin{axiom}[Field equation]\label{ax:FE}
On physical configurations $T=0$, with the memory channel oriented by Definition~\ref{def:orient}:
\begin{equation}\label{eq:FE}
\alpha J(x)=(1-\alpha)\,e^{-i\sigma\alpha\pi/2}\int_{J^-(x)}k(\tau;x,x')\,\mathcal P(x,x')\,J(x')\,d^4x' .
\end{equation}
\end{axiom}

\begin{remark}\label{rem:status}
Axiom~\ref{ax:FE} replaces Posit~1 of~\cite{Skeleton}. It assumes no unitarity, no generator, and no identification $\hat H=f_{\rm map}[\hat T]$.

The alternative balance law $\partial_tT=0$ gives the same bounded admissible solutions. If $J$ is admissible then so is $T$, because $\hat T$ is a Fourier multiplier of $\hat J$. A constant admissible $T$ has no support at zero frequency, so it vanishes.
\end{remark}

%======================================================================
\section{Unconditional lemmas}

\begin{lemma}[Weyl convergence]\label{lem:weyl}
Let $\omega\in\mathbb R\setminus\{0\}$. The improper integral $\int_0^\infty s^{\alpha-1}e^{i\omega s}\,ds$ converges if and only if $0<\alpha<1$. When it converges, its value is $\Gamma(\alpha)(-i\omega)^{-\alpha}$. Consequently $W^\alpha e^{-iEt/\hbar}=(-iE/\hbar)^{-\alpha}e^{-iEt/\hbar}$ for $0<\alpha<1$. At $\alpha=1$ the right-hand side exists only as the limit $\alpha\to1^-$.
\end{lemma}

\begin{proof}
\emph{Convergence for $0<\alpha<1$.} Near $s=0$ the integrand is integrable. On $[1,\infty)$, $s^{\alpha-1}$ decreases monotonically to $0$ and $\int_1^Xe^{i\omega s}ds$ is bounded, so Dirichlet's test applies.

\emph{Divergence for $\alpha\ge1$.} If $\alpha=1$, the partial integrals $(e^{i\omega X}-1)/(i\omega)$ have no limit. If $\alpha>1$, the amplitude $s^{\alpha-1}$ does not tend to zero, so again there is no limit.

\emph{Value.} Rotate the contour onto the ray $\arg s=\tfrac\pi2\sgn\omega$, and use Jordan's lemma for the arc. \end{proof}

\begin{lemma}[Inversion]\label{lem:inv}
On the sector, $\mathcal P$ is time-translation invariant and $\hat R_\Hil$ commutes with $W^\alpha$. On bounded admissible $J$, \eqref{eq:FE} is equivalent to
\begin{equation}\label{eq:FEinv}
\alpha\,(\sigma i\tau_c)^{\alpha}\,D^\alpha J=\Lambda J .
\end{equation}
\end{lemma}

\begin{proof}
On the sector, \eqref{eq:FE} reads $\alpha J=\eps\,e^{-i\sigma\alpha\pi/2}\tau_c^{-\alpha}\hat R_\Hil W^\alpha J$. Apply $D^\alpha$ and multiply by $(\sigma i\tau_c)^\alpha=e^{i\sigma\alpha\pi/2}\tau_c^\alpha$. For the converse, apply $W^\alpha$. Both maps are injective on admissible functions. \end{proof}

\begin{lemma}[Phase cancellation]\label{lem:phase}
Let $g_\alpha(E)\equiv\alpha(\sigma i)^\alpha(-iE/E_c)^\alpha$ be the symbol of the left side of~\eqref{eq:FEinv}. For $E\neq0$,
\[
g_\alpha(E)=\alpha\lvert E/E_c\rvert^\alpha\,e^{\,i\alpha\pi(\sigma-\sgn E)/2}.
\]
For $\alpha\in(0,1)\cup(1,2)$, $g_\alpha(E)$ is real if and only if $\sgn E=\sigma$, and then it is positive and strictly increasing in $\lvert E\rvert$. At $\alpha=1$, $g_1(E)=\sigma E/E_c$ is real for both signs of $E$. Finally, setting $f_\alpha(\lambda)\equiv\sigma E_c(\lambda/\alpha)^{1/\alpha}$ for $\lambda\ge0$ gives $g_\alpha(f_\alpha(\lambda))=\lambda$.
\end{lemma}

\begin{proof}
The principal arguments are $\arg(\sigma i)=\sigma\pi/2$ and $\arg(-iE)=-\sgn(E)\,\pi/2$. Hence $\arg g_\alpha=\alpha\pi(\sigma-\sgn E)/2\in\{0,\pm\alpha\pi\}$, and $\pm\alpha\pi\notin\pi\mathbb Z$ for $\alpha\notin\{0,1,2\}$. The inverse relation follows by direct substitution. \end{proof}

\begin{lemma}[Operator facts]\label{lem:ops}
\begin{enumerate}
\item $\hat R_\Hil=r_0\nabla^2+r_1$ is self-adjoint on $H^2(\mathbb R^3)$: it is unitarily equivalent to multiplication by $r_1-r_0\lvert k\rvert^2$. A bounded $x$-dependent $r_1$ preserves self-adjointness, by Kato--Rellich~\cite{Kato}.
\item $\lVert\hat R_\Hil\rVert=\infty$, so the spectral inclusion in~\cite[Lemma~4.4]{Skeleton} is vacuous.
\item Exchanging the kernel limit with $\Pi_\Hil$ requires strong-resolvent convergence of the projected operators~\cite{RS1}. The bound $\lVert\Pi_\Hil\rVert=1$ alone does not suffice.
\end{enumerate}
\end{lemma}

%======================================================================
\section{Derived unitarity and the collapse}

Write $\Hil_+\equiv\Ran\one_{(0,\infty)}(\Lambda)$.

\begin{theorem}[Bounded solutions are unitary]\label{thm:eternal}
Fix $\alpha\in(0,1)$ and $\sigma$. Let $J\in C_b(\mathbb R;\Hil)$ be admissible and solve~\eqref{eq:FEinv} in the sense of $\Hil$-valued tempered distributions. Then:
\begin{enumerate}
\item[(i)] $J(t)=e^{-itf_\alpha(\Lambda)/\hbar}\psi$ for some $\psi\in\Hil_+$. Conversely, every such $J$ is a solution.
\item[(ii)] $U_\alpha(t)\equiv e^{-itf_\alpha(\Lambda)/\hbar}\big|_{\Hil_+}$ is a strongly continuous one-parameter unitary group with self-adjoint generator $f_\alpha(\Lambda)$. In particular $\lVert J(t)\rVert=\lVert J(0)\rVert$.
\item[(iii)] Every frequency present satisfies $\sigma E>0$, and $\Ran\one_{(-\infty,0)}(\Lambda)$ carries no solution.
\end{enumerate}
\end{theorem}

\begin{proof}
\emph{Step 1: fiber equation.} Fourier-transform in $t$ and pass to the spectral representation of $\Lambda$. Fiberwise, \eqref{eq:FEinv} becomes $\big(g_\alpha(E)-\lambda\big)\hat J_\lambda(E)=0$.

\emph{Step 2: zeros.} By Lemma~\ref{lem:phase}:
\begin{itemize}
\item for $\lambda>0$, the continuous function $g_\alpha-\lambda$ vanishes only at the simple zero $E_\lambda=f_\alpha(\lambda)$, where $\sigma E_\lambda>0$;
\item for $\lambda<0$, it never vanishes;
\item for $\lambda=0$, its only zero is $E=0$, which admissibility excludes.
\end{itemize}

\emph{Step 3: support.} A tempered distribution annihilated by a continuous function is supported on that function's zero set. So $\hat J_\lambda$ is a finite sum of $\delta^{(n)}(E-E_\lambda)$, and $\hat J_\lambda=0$ when $\lambda\le0$. The terms with $n\ge1$ produce growth $t^n$ and are excluded by boundedness. Hence $J_\lambda(t)=c_\lambda e^{-iE_\lambda t/\hbar}$, which is~(i).

\emph{Step 4: converse.} By $g_\alpha(f_\alpha(\lambda))=\lambda$, each such $J$ satisfies the fiber equation. Its spectrum lies in $\sigma E>0$, so it is admissible.

\emph{Step 5: unitarity.} Functional calculus makes $U_\alpha$ a strongly continuous unitary group. Stone's theorem~\cite{Stone,RS1} identifies its generator as $f_\alpha(\Lambda)$, which is self-adjoint because $f_\alpha$ is real and Borel. Statement~(iii) is Step~2. \end{proof}

\begin{remark}
Theorem~\ref{thm:eternal} is the derived replacement for Posit~1 of~\cite{Skeleton}. The fiberwise argument of Steps~1--3 is stated at the level of rigor of the companion papers. The measurability details are standard but are not written out here.
\end{remark}

\begin{theorem}[Collapse to Schr\"odinger]\label{thm:collapse}
Take the double scaling with $\Lambda$ fixed. For every $\psi\in\Hil_+$ and every $T>0$,
\[
\sup_{\lvert t\rvert\le T}\big\lVert U_\alpha(t)\psi-e^{-itH/\hbar}\psi\big\rVert\xrightarrow[\alpha\to1^-]{}0,
\qquad H\equiv\sigma E_c\Lambda .
\]
If $\psi$ has compact spectral support in $(0,\infty)$, then
\begin{equation}\label{eq:rate}
\big\lVert U_\alpha(t)\psi-e^{-itH/\hbar}\psi\big\rVert\le\frac{\lvert t\rvert}{\hbar}\big\lVert(f_\alpha-f_1)(\Lambda)\psi\big\rVert,
\qquad
(f_\alpha-f_1)(\lambda)=\sigma E_c\,\eps\,\lambda(1+\ln\lambda)+O(\eps^2).
\end{equation}
The limiting equation is $i\hbar\,\partial_tJ=HJ$. On the spinless, gauge-neutral sector,
\[
H=-\frac{\hbar^2}{2m}\nabla^2+V_0,\qquad \frac{\hbar^2}{2m}=-\sigma\eps E_cr_0,\qquad V_0=\sigma\eps E_cr_1 .
\]
Positive mass with $r_0>0$ forces $\sigma=-1$. Then $\hbar^2/2m=\eps E_cr_0$, which is the mass formula of~\cite[Thm.~6.1(iii)]{Skeleton} with its units restored, and $\Hil_+=\Ran\one_{(-\infty,0)}(H)$.
\end{theorem}

\begin{proof}
\emph{Pointwise limit.} Since $f_\alpha\to f_1=\sigma E_c\lambda$ pointwise on $[0,\infty)$, the integrand $\lvert e^{-itf_\alpha(\lambda)/\hbar}-e^{-itf_1(\lambda)/\hbar}\rvert^2\le4$ tends to zero. Dominated convergence in the spectral measure of $\psi$ gives the strong limit. Uniformity for $\lvert t\rvert\le T$ follows from $\lvert e^{ia}-e^{ib}\rvert\le\min(2,\lvert a-b\rvert)$ together with the local uniformity of $f_\alpha\to f_1$.

\emph{Rate.} The same inequality gives the bound in~\eqref{eq:rate}. Expanding $\ln\!\big((\lambda/\alpha)^{1/\alpha}\big)=(1+\eps+\dots)(\ln\lambda+\eps+\dots)$ gives the stated difference.

\emph{Identification.} $f_1(\Lambda)=\sigma E_c\Lambda$, and on the sector $\Lambda$ acts as multiplication by $\eps(r_1-r_0\lvert k\rvert^2)$. \end{proof}

\begin{corollary}[The skeleton and the edge are two scalings of one limit]\label{cor:scalings}
\begin{enumerate}
\item Hold $\hat R_\Hil$ and $\tau_c$ fixed, as in~\cite[Lemma~4.1]{Skeleton}. The memory rate vanishes (Proposition~\ref{prop:rate}(1)) and $f_\alpha(\eps\hat R_\Hil)\to0$. Then $U_\alpha(t)\to\one$ strongly, and $m=\hbar^2/(2\eps E_cr_0)\to\infty$. The carrier freezes. This is the kernel-blind skeleton of~\cite[Thm.~5.1]{Skeleton}, up to its unobservable global phase.
\item Hold the memory rate fixed, equivalently $\Lambda$. Then Theorem~\ref{thm:collapse} applies: the memory stays live and becomes $H$.
\end{enumerate}
The ``corner'' and the ``edge'' of~\cite[Rem.~5.2, Table~7.4]{Skeleton} are therefore one limit taken in two ways.
\end{corollary}

\begin{proposition}[Response to a source]\label{prop:switch}
Let $\alpha\in(\tfrac23,1)$. Let $S$ be a bounded source with zero past, and consider $T=S$, that is, \eqref{eq:FE} with $S$ on the left. Its causal solution is $\hat J(E)=\hat G(E)\hat S(E)$, with
\[
\hat G(E)=\frac{(\sigma i)^\alpha(-iE/E_c)^\alpha}{g_\alpha(E)-\Lambda},
\]
evaluated just above the real axis. On each spectral value $\lambda\neq0$, $\hat G$ is analytic in ${\rm Im}\,E>0$; its other singularities are zeros of $g_\alpha-\lambda$ and the branch cut on the negative imaginary axis.
\begin{enumerate}
\item For $\lambda>0$ the zero is the real point $E_\lambda=f_\alpha(\lambda)$. After the source has acted, the response is the stationary solution of Theorem~\ref{thm:eternal}, with amplitude
\[
\Big\lvert\frac{(\lambda/\alpha)\,\hat S(E_\lambda)}{g_\alpha'(E_\lambda)}\Big\rvert ,
\]
plus a tail from the branch cut that decays as a power of $t$.
\item For $\lambda<0$ the zero is at ${\rm Im}\,E=-\pi\eps\lvert E\rvert\big(1+O(\eps)\big)<0$. Such a mode exists only as a causal transient that decays at this rate. It has no bounded solution on the whole time line, consistent with Theorem~\ref{thm:eternal}(iii).
\end{enumerate}
\end{proposition}

\begin{proof}
\emph{Response function.} Fourier-transform $\alpha J-\eps e^{-i\sigma\alpha\pi/2}\tau_c^{-\alpha}\hat R_\Hil W^\alpha J=S$ and multiply by $(\sigma i)^\alpha(-iE/E_c)^\alpha$. This gives $\hat G$.

\emph{Analyticity.} The factor $(-iE)^\alpha$, principal branch, is analytic away from $-iE\in(-\infty,0]$, which is the negative imaginary $E$-axis. This analyticity in the upper half-plane is the causality of $W^\alpha$.

\emph{Location of the zeros.} Put $w=-iE/E_c$, so that ${\rm Im}\,E=E_c\,{\rm Re}\,w$, and $g_\alpha=\lambda$ becomes $w^\alpha=(\lambda/\alpha)(\sigma i)^{-\alpha}$.
\begin{itemize}
\item For $\lambda>0$: $w=-i\sigma(\lambda/\alpha)^{1/\alpha}$, so ${\rm Re}\,w=0$ and the zero is real.
\item For $\lambda<0$: the principal root has $\arg w=\pm(\pi/\alpha-\pi/2)$. A principal root exists iff $\pi-\alpha\pi/2<\alpha\pi$, that is, $\alpha>\tfrac23$. Then ${\rm Re}\,w=\lvert w\rvert\sin(\pi/\alpha)$, and for $\alpha=1-\eps$ this is $-\lvert w\rvert\sin(\pi\eps)\,(1+O(\eps))$.
\end{itemize}

\emph{Late-time response.} For $t$ after the source, close the inversion contour in the lower half-plane. The residue at the real zero gives (1). The cut gives the power-law tail~\cite{GKMR,Podlubny}. \end{proof}

\begin{remark}
A Caputo initial-value problem with $J(0)=J_0$ corresponds to the frozen past $J\equiv J_0$ for $t<0$. Its stationary amplitude is $\alpha^{-1}\lvert J_0\rvert$, not $\lvert J_0\rvert$. That factor belongs to the frozen past. The frozen past is not admissible for Axiom~\ref{ax:FE}, so the factor is not a prediction of the field equation.
\end{remark}

\begin{theorem}[The limit is one-sided]\label{thm:onesided}
\begin{enumerate}
\item \emph{Convergence.} $W^\alpha$ converges on oscillating carriers only for $\alpha<1$ (Lemma~\ref{lem:weyl}). The memory term of~\eqref{eq:FE} at $\alpha=1$ exists only as the limit from below.
\item \emph{Sign.} With $\sigma$ fixed, the memory-generated coefficient $\hbar^2/2m=-\sigma\eps E_cr_0$ is odd in the memory rate $\eps E_c/\hbar$. Positive mass can hold on one side of $\alpha=1$ only.
\item \emph{Runaway or preacceleration.} For $\alpha\in(1,2)$, take the source problem with zero past. There the Riemann--Liouville integral is finite, and $\hat G$ is defined by the continued symbol. On $\Ran\one_{(-\infty,0)}(\Lambda)$, $g_\alpha-\lambda$ has a zero with ${\rm Im}\,E\approx+\pi\lvert\eps\rvert\lvert E\rvert>0$, in the upper half-plane. Then:
\begin{itemize}
\item the causal response, obtained by inverting above every singularity, grows at that rate;
\item the bounded response, obtained by inverting on the real axis, is nonzero before the source acts.
\end{itemize}
This is the Abraham--Lorentz--Dirac dichotomy of the third-derivative panel of~\cite{CAL}, runaway or preacceleration, appearing on the wrong side of $\alpha=1$.
\end{enumerate}
\end{theorem}

\begin{proof}
Part~(1) is Lemma~\ref{lem:weyl}. Part~(2) is immediate from Theorem~\ref{thm:collapse}. For part~(3), the argument in the proof of Proposition~\ref{prop:switch} gives ${\rm Re}\,w=\lvert w\rvert\sin(\pi/\alpha)$, which is positive for $\alpha\in(1,2)$. A zero of the response function in the upper half-plane rules out a response that is both causal and bounded. \end{proof}

%======================================================================
\section{The two Markovian regimes, causal structure, and interaction}\label{sec:regimes}

The framework has two Markovian regimes. One is the corner $\alpha=1$, where the past is disconnected. The other is the $\delta$ (Dirac) limit of the kernel, where only the immediate past survives. This section shows that neither yields unitary Schr\"odinger evolution on its own, and identifies what the edge requires instead. It then records two structural limits of the edge: it is Galilean, and it carries no interaction.

\subsection{The corner: the past disconnects}

\begin{proposition}[Memory rate]\label{prop:rate}
$\rho_\eps(\tau)=\eps\,\tau_c^{-1+\eps}\tau^{-\eps}/\Gamma(1-\eps)$.
\begin{enumerate}
\item With $\tau_c$ fixed, $\rho_\eps\to0$ uniformly on compact lag intervals. The past disconnects even from below, and the edge degenerates to the skeleton (Corollary~\ref{cor:scalings}(1)).
\item With $\tau_c=\eps/\kappa$, $\rho_\eps(\tau)\to\kappa$ for every fixed $\tau>0$. The kernel becomes the flat accumulator of rate $\kappa$, not a $\delta$.
\end{enumerate}
\end{proposition}

\begin{proof}
For~(1), $\rho_\eps\le C\eps\,\tau^{-\eps}$ on compacts. For~(2), $\rho_\eps=\kappa^{1-\eps}\eps^{\eps}\tau^{-\eps}/\Gamma(1-\eps)$, and $\eps^\eps\to1$, $\tau^{-\eps}\to1$, $\Gamma(1-\eps)\to1$. \end{proof}

\subsection{The $\delta$ limit: no stationary band}

\begin{theorem}[Normalizable kernels have no stationary band]\label{thm:normalizable}
Replace the power law in~\eqref{eq:FE} by a lag weight $w\ge0$ with $\int_0^\infty w=1$, of any width. This includes every family $w\to\delta$.

On a spectral value $\lambda$ of $\hat R_\Hil$, a plane wave $e^{-iEt/\hbar}$ with real $E$ solves the field equation only if $\eps\lvert\lambda\rvert\ge\alpha$. Hence on $\{\eps\lvert\lambda\rvert<\alpha\}$ the field equation has no bounded solution on the whole time line. That subspace contains every small value of $\Lambda=\eps\lambda$, and so the entire low-energy band.
\end{theorem}

\begin{proof}
On a plane wave the field equation reads $\alpha=\eps\,e^{-i\sigma\theta}\lambda\,\hat w(E)$, where $\hat w(E)=\int w(\tau)e^{iE\tau/\hbar}d\tau$. For real $E$, $\lvert\hat w(E)\rvert\le\int w=1$, so a real root needs $\alpha\le\eps\lvert\lambda\rvert$. The support argument of Theorem~\ref{thm:eternal} (Steps~1--3) then excludes bounded solutions wherever no real root exists. \end{proof}

\begin{example}[A narrow exponential kernel]\label{ex:exp}
Take $w(\tau)=e^{-\tau/\tau_w}/\tau_w$, so that $\hat w(E)=1/(1-iE\tau_w/\hbar)$. The root is exact:
\[
E=-\frac{i\hbar}{\tau_w}\Big[1-\frac{\eps\lambda}{\alpha}e^{-i\sigma\theta}\Big].
\]
Hold the would-be generator $\kappa\equiv\eps\lambda\hbar/\tau_w$ fixed. Then
\[
{\rm Re}\,E=\sigma\frac{\kappa}{\alpha}\sin\theta\to\sigma\kappa,
\qquad
{\rm Im}\,E=-\frac{\hbar}{\tau_w}+\frac{\kappa}{\alpha}\cos\theta .
\]
The real part is the Schr\"odinger frequency. But every mode decays at the inverse kernel width, which becomes infinite as $\tau_w\to0$ (Table~\ref{tab:exp}). A narrow kernel's symbol is nearly real, so it cannot absorb the channel factor $e^{-i\sigma\theta}\to\mp i$, which therefore appears as damping rather than as a phase.
\end{example}

\begin{remark}[Why the power law]\label{rem:powerlaw}
The power-law kernel is not normalizable, since $\int_0^\infty\tau^{\alpha-1}d\tau=\infty$. Its symbol $(-iE\tau_c/\hbar)^{-\alpha}$ is unbounded as $E\to0$. So $\lvert\text{symbol}\rvert=\alpha/(\eps\lambda)$ is met at a finite real energy $E=E_c(\eps\lambda/\alpha)^{1/\alpha}$, which is the stationary band of Theorem~\ref{thm:eternal}. The symbol's phase is exactly the channel angle $\alpha\pi/2$, which is how~\cite[Def.~2.2]{Skeleton} defines the channel angle.

Infinite total memory weight, multiplied by a vanishing $\eps$, gives the finite rate of Proposition~\ref{prop:rate}(2). The edge is therefore a third regime, distinct from both Markovian ones: a law with memory (power law, full width) whose flow is Markovian (Theorem~\ref{thm:eternal}). In the table of~\cite{Skeleton}, the Schr\"odinger form belongs in the full-width (power-law) column at $1^-$, not in the $\delta$ column. Posit~1 was covering the $\delta$ column's lack of a stationary band.
\end{remark}

\subsection{Causal structure: the edge is Galilean}\label{sec:galilean}

In the framework the Lorentz signature appears with the quadratic closure $JJ^\dagger$ at $\alpha=2$, through the invariant $\det(JJ^\dagger)=\lvert\det J\rvert^2$~\cite{CAL}. The same closure gives the Born readout through the trace, $\tfrac12{\rm tr}(JJ^\dagger)=\sum_a\lvert Z_a\rvert^2$. That trace is a density, the time component of a current; it is not a Lorentz scalar. Below the closure there is no determinant, so there is no metric signature and no light cone.

At the $\alpha\to1^-$ edge the causal structure is therefore time ordering alone: $J^-(x)$ reduces to the time-ordered past $\{t'<t\}$, with no bound on spatial reach. This is the Galilean structure of the Schr\"odinger equation, whose free propagator is nonzero at every distance for every $t>0$. Three things follow.
\begin{enumerate}
\item The lag-independent spatial profile of Definition~\ref{def:kernel} is admissible.
\item Instantaneous spatial reach in a $\delta$ route is admissible as well. That route fails on unitarity (Theorem~\ref{thm:normalizable}), not on causality.
\item A light cone imposed at the edge imports $\alpha=2$ structure. For reference only, Appendix~\ref{app:cone} records what imposing one would give.
\end{enumerate}

\subsection{The edge carries no interaction}\label{sec:nointeraction}

\begin{remark}[Single legs]\label{rem:nointeraction}
In~\eqref{eq:FE} the past enters only through transports $\mathcal P(x,x')$ from a point of the past to the present. The generator $H=\sigma E_c\Lambda$ of Theorem~\ref{thm:collapse} is assembled from these open legs, and no term transports the present back. Hence:
\begin{enumerate}
\item there is no reciprocal, two-leg closure, so there is no interaction in the sense of~\cite{CAL}: no back-reaction and no two-body coupling;
\item the carrier generates no holonomy, since holonomy needs a closed path~\cite{CAL};
\item what the edge can carry is a fixed background acting one way, through a bounded $r_1(x)$ (Lemma~\ref{lem:ops}) or a prescribed transport. That is a test particle that does not act back.
\end{enumerate}
Two-body interaction and entangling dynamics require the composite carrier $J_{AB}\in M_2(\mathbb C)\otimes M_2(\mathbb C)\cong M_4(\mathbb C)$ and its closure $J_{AB}J_{AB}^\dagger$ at $\alpha=2$. This agrees with the restriction of~\cite[Thm.~6.1]{Skeleton} to the free form: that restriction is structural, not a matter of convenience.
\end{remark}

%======================================================================
\section{Numerical verification}

All statements above were run with the script \texttt{scripts/cal\_schrodinger\_collapse.py} in the repository \url{https://github.com/JewJitsu11B/Liorhybrid}. It uses NumPy only and runs in about ten seconds. Units are $\hbar=\tau_c=1$, so $E_c=1$. Methods are described in Appendix~\ref{app:num}.

\emph{Bounded solutions}, with $\sigma=-1$, $m=1$, and cap $E_{\rm cap}=4$. The test state is a Gaussian packet with $k_0=1$ and width $2$; its weight above the cap is $10^{-13}$.

\begin{table}[h]
\centering\small
\begin{tabular}{cccccc}
\toprule
$\eps$ & $\max\lvert E_\alpha-E_1\rvert$ & $\eps\max\lambda\lvert1+\ln\lambda\rvert$ & ratio & phase-aligned error, $t\le40$ & norm drift\\
\midrule
$10^{-1}$ & $1.246$ & $0.954$ & $1.305$ & $1.41$ & $3\times10^{-16}$\\
$10^{-2}$ & $9.78\times10^{-2}$ & $9.55\times10^{-2}$ & $1.025$ & $3.30\times10^{-1}$ & $3\times10^{-16}$\\
$10^{-3}$ & $9.57\times10^{-3}$ & $9.55\times10^{-3}$ & $1.002$ & $3.27\times10^{-2}$ & $3\times10^{-16}$\\
$3\times10^{-4}$ & $2.87\times10^{-3}$ & $2.86\times10^{-3}$ & $1.001$ & $9.80\times10^{-3}$ & $3\times10^{-16}$\\
\bottomrule
\end{tabular}
\caption{Theorems~\ref{thm:eternal}--\ref{thm:collapse}. The energy error matches the first-order term of~\eqref{eq:rate}, and the norm is conserved to machine precision. The raw packet error is dominated by a global phase drift from the constant term. With that phase removed, the error is linear in $\eps$, at about $33\,\eps$ over $t\le40$.}
\end{table}

\emph{Source-driven runs.} The integral equation $T=S$ is time-stepped directly, with Weyl integration over a zero past and the full history kept. The source is a Gaussian pulse at $t_0=10$ of width $2$, tuned to the mode. No initial-value data enter.

\begin{table}[h]
\centering\small
\begin{tabular}{ccccc}
\toprule
$\alpha$ & $h$ & mean $\lvert J\rvert$, $t\in[80,120]$ & residue prediction & ratio\\
\midrule
$0.900$ & $0.02$ & $6.95236$ & $6.95787$ & $0.99921$\\
$0.900$ & $0.01$ & $6.95674$ & $6.95787$ & $0.99984$\\
$0.990$ & $0.01$ & $5.16659$ & $5.16724$ & $0.99987$\\
$0.999$ & $0.01$ & $5.02771$ & $5.02833$ & $0.99988$\\
$1.000$ & $0.01$ & $5.01264$ & $5.01326$ & $0.99988$\\
$1.010$ & $0.01$ & $4.86571$ & $4.86630$ & $0.99988$\\
\bottomrule
\end{tabular}
\caption{Below the cap ($\Lambda=+1$), Proposition~\ref{prop:switch}(1). After the pulse the mode oscillates forever at the amplitude of the residue at the real zero. This is the stationary solution of Theorem~\ref{thm:eternal}, excited causally. The ratio converges as $h$ is halved.}
\end{table}

\begin{table}[h]
\centering\small
\begin{tabular}{ccccc}
\toprule
$\alpha$ & zero of $g_\alpha-\Lambda$ & $-{\rm Im}\,E$ & $\pi\eps\lvert E\rvert$ & measured decay rate\\
\midrule
$0.970$ & $1.027032-0.100104\,i$ & $+1.0010\times10^{-1}$ & $+9.73\times10^{-2}$ & $+1.0007\times10^{-1}$\\
$0.990$ & $1.009695-0.032052\,i$ & $+3.2052\times10^{-2}$ & $+3.17\times10^{-2}$ & $+3.2049\times10^{-2}$\\
$0.999$ & $1.000997-0.003148\,i$ & $+3.1479\times10^{-3}$ & $+3.14\times10^{-3}$ & $+3.1478\times10^{-3}$\\
$1.000$ & $1.000000$ & $0$ & $0$ & $+2.5\times10^{-7}$\\
$1.001$ & $0.998997+0.003135\,i$ & $-3.1353\times10^{-3}$ & $-3.14\times10^{-3}$ & $-3.1348\times10^{-3}$\\
$1.010$ & $0.989718+0.030795\,i$ & $-3.0795\times10^{-2}$ & $-3.11\times10^{-2}$ & $-3.0792\times10^{-2}$\\
$1.030$ & $0.967645+0.088790\,i$ & $-8.8790\times10^{-2}$ & $-9.16\times10^{-2}$ & $-8.8782\times10^{-2}$\\
\bottomrule
\end{tabular}
\caption{Above the cap ($\Lambda=-1$), Proposition~\ref{prop:switch}(2) and Theorem~\ref{thm:onesided}(3). From below the zero lies in ${\rm Im}\,E<0$, giving causal decay. At $\alpha=1$ it lies on the real axis. From above it lies in ${\rm Im}\,E>0$, and the causal response grows. The measured rate is $(\ln\lvert J(40)\rvert-\ln\lvert J(80)\rvert)/40$.}
\end{table}

\begin{table}[h]
\centering\small
\begin{tabular}{cccc}
\toprule
$\alpha$ & $\lvert J(t_0-10)\rvert$ & $\lvert J(t_0-20)\rvert$ & $\lvert J(t_0+40)\rvert$\\
\midrule
$0.99$ & $4.0\times10^{-6}$ & $7.5\times10^{-11}$ & $1.44$\\
$1.01$ & $3.58$ & $2.63$ & $4.2\times10^{-6}$\\
\bottomrule
\end{tabular}
\caption{Bounded responses, from inversion on the real axis, for $\Lambda=-1$. From below the bounded response is causal: what appears before the pulse is only the Gaussian's own tail. From above, keeping the response bounded makes it precede its source. That is preacceleration, Theorem~\ref{thm:onesided}(3).}
\end{table}

\emph{The two Markovian regimes.}

\begin{table}[h]
\centering\small
\begin{tabular}{ccccccc}
\toprule
 & \multicolumn{2}{c}{$\tau_c=1$ fixed} & & \multicolumn{3}{c}{$\tau_c=\eps$}\\
\cmidrule{2-3}\cmidrule{5-7}
$\eps$ & $\tau=1$ & $\tau=100$ & & $\tau=0.01$ & $\tau=1$ & $\tau=100$\\
\midrule
$10^{-1}$ & $9.36\times10^{-2}$ & $5.90\times10^{-2}$ & & $1.178$ & $0.743$ & $0.469$\\
$10^{-2}$ & $9.94\times10^{-3}$ & $9.49\times10^{-3}$ & & $0.994$ & $0.949$ & $0.907$\\
$10^{-3}$ & $9.99\times10^{-4}$ & $9.95\times10^{-4}$ & & $0.997$ & $0.993$ & $0.988$\\
$10^{-4}$ & $1.00\times10^{-4}$ & $1.00\times10^{-4}$ & & $1.000$ & $0.999$ & $0.999$\\
\bottomrule
\end{tabular}
\caption{Memory rate $\rho_\eps(\tau)$, Proposition~\ref{prop:rate}. With $\tau_c$ fixed the rate is about $\eps$, and the past disconnects even from below. With $\tau_c\propto\eps$ the rate tends to $1$ at every lag: a flat accumulator, not a $\delta$.}
\end{table}

\begin{table}[h]
\centering\small
\begin{tabular}{ccccc}
\toprule
$\eps$ & $\tau_w$ & ${\rm Re}\,E$ & ${\rm Im}\,E$ & $-1/\tau_w$\\
\midrule
$10^{-2}$ & $1$ & $+1.0100$ & $-0.984$ & $-1$\\
$10^{-2}$ & $0.01$ & $+1.0100$ & $-99.98$ & $-100$\\
$10^{-3}$ & $1$ & $+1.0010$ & $-0.998$ & $-1$\\
$10^{-3}$ & $0.1$ & $+1.0010$ & $-9.998$ & $-10$\\
$10^{-3}$ & $0.001$ & $+1.0010$ & $-1000.0$ & $-1000$\\
\bottomrule
\end{tabular}
\caption{The normalizable exponential kernel of Example~\ref{ex:exp}, with generator $\kappa=1$ and $\sigma=+1$. The real part reproduces the Schr\"odinger frequency, but every mode decays at about $1/\tau_w$. The power law with the same $\kappa$ gives the real root $1.0010+0i$ at $\eps=10^{-3}$.}
\label{tab:exp}
\end{table}

\clearpage
%======================================================================
\section{Remarks on scope}

\begin{remark}[Status ledger]
\emph{Inputs:}
\begin{itemize}
\item the field equation (Axiom~\ref{ax:FE});
\item the power-law kernel class (Definition~\ref{def:kernel});
\item time ordering as the causal structure below the $\alpha=2$ closure;
\item the double scaling at fixed memory rate (Definition~\ref{def:double}).
\end{itemize}
\emph{Derived from these:} every lemma, theorem, proposition and corollary of \S\S4--6.

\emph{Not claimed:}
\begin{itemize}
\item a value for the mass scale;
\item interaction, entanglement, Lorentz structure, or a gauge field strength generated by the carrier, all of which belong to $\alpha=2$;
\item any initial-value (Caputo) statement, since a frozen past is not admissible.
\end{itemize}
\end{remark}

\begin{remark}[What changes relative to~\cite{Skeleton}]
\begin{itemize}
\item Posit~1, and Theorem~6.1(i) built on it, are replaced by Axiom~\ref{ax:FE} together with Theorems~\ref{thm:eternal} and~\ref{thm:collapse}.
\item The $\delta$ column at the $1^-$ edge has no stationary band once the channel phase is used (Theorem~\ref{thm:normalizable}). That is the gap Posit~1 was covering. The Schr\"odinger form belongs in the power-law (full-width) column.
\item Below $\alpha=2$ the causal structure is Galilean, so the $\delta$ route's instantaneous spatial reach is admissible. Its obstruction is unitarity, not causality.
\item The free form of Theorem~6.1(ii) is structural: the edge carries no interaction (Remark~\ref{rem:nointeraction}).
\item Axiom~1 should require $k$ and $\mathcal P$ to be real, not $J$.
\item In Lemma~4.3, real plus symmetric gives a symmetric operator; self-adjointness is Lemma~\ref{lem:ops}(1).
\item In Lemma~4.4, $\lVert\hat R\rVert=\infty$, and exchanging the limit with the projection needs strong-resolvent convergence.
\item The mass formula of Theorem~6.1(iii) holds at fixed memory rate. With $\tau_c$ fixed the mass diverges (Corollary~\ref{cor:scalings}(1)).
\item The constant shift in Eq.~(11) is inert only at $\alpha=1$. For $\alpha<1$ it sets the stationary sector, by Theorem~\ref{thm:eternal}(iii).
\item The $1^-$ edge of Table~7.4 is one-sided by Theorem~\ref{thm:onesided}, not by convention.
\end{itemize}
\end{remark}

\begin{remark}[Dynamics is memory]
Corollary~\ref{cor:scalings} and Theorem~\ref{thm:normalizable} make~\cite[Rem.~7.3]{Skeleton} precise.
\begin{itemize}
\item If the memory weight vanishes at fixed time unit, the past disconnects and the carrier freezes.
\item If it vanishes at fixed memory rate, the memory becomes the Hamiltonian, but only if the memory has infinite total weight. A normalizable memory, however narrow, cannot supply a unitary generator.
\end{itemize}
The Schr\"odinger generator is the image of an infinite accumulated past under the double scaling, transported one way, with no return leg.
\end{remark}

\appendix
%======================================================================
\section{Notation}

\begin{tabular}{ll}
$\alpha$, $\eps=1-\alpha$ & fractional order and memory weight\\
$\sigma$, $\theta=\alpha\pi/2$ & channel orientation and channel angle\\
$W^\alpha$, $D^\alpha$ & causal Weyl integral and its inverse (Definition~\ref{def:weyl})\\
$\tau_c$, $E_c=\hbar/\tau_c$ & memory time unit and its energy\\
$\rho_\eps(\tau)$ & memory rate at lag $\tau$ (Definition~\ref{def:kernel})\\
$g$, $\sigma_g^2$ & spatial profile and its second moment\\
$\hat R_\Hil=r_0\nabla^2+r_1$ & projected memory operator\\
$\Lambda=\eps\hat R_\Hil$ & double-scaled memory operator\\
$g_\alpha$, $f_\alpha$ & symbol of the left side of~\eqref{eq:FEinv}, and stationary energy\\
$\Hil_+$ & $\Ran\one_{(0,\infty)}(\Lambda)$, the stationary sector\\
$w$, $\hat w$, $\tau_w$ & normalizable lag weight, its symbol, its width\\
\end{tabular}

\section{Numerical methods}\label{app:num}

\emph{Time stepping.} The Weyl integral over a zero past is discretized by BDF2 convolution quadrature~\cite{Lubich}, with weights $q_j$ generated by $\big((1-\zeta)+(1-\zeta)^2/2\big)^{-\alpha}$. The full history is kept. Each step solves
\[
(\alpha-\kappa h^\alpha q_0)\,J_n=S_n+\kappa h^\alpha\sum_{j\ge1}q_jJ_{n-j},\qquad \kappa=\Lambda e^{-i\sigma\alpha\pi/2}.
\]
This is the integral field equation itself. No derivative and no initial-value data are used.

\emph{Zeros and residues.} Zeros of $g_\alpha-\Lambda$ are found by Newton iteration. Residue amplitudes use $\hat S(E)=\sqrt{2\pi}\,w\,e^{-w^2(E-\Omega)^2/2}e^{i(E-\Omega)t_0}$. The bounded responses come from direct inversion on the real axis, on a grid of $1.6\times10^5$ points over $\lvert E\rvert\le40$.

\emph{Packets.} Packets use a periodic grid of $2048$ points on a box of length $400$, with a unitary FFT. Modes with amplitude below $10^{-10}$ of the maximum are dropped.

\emph{Markovian regimes.} The memory rate is evaluated in closed form. The exponential-kernel roots are the closed form of Example~\ref{ex:exp}.

%======================================================================
\section{For reference: a light cone imposed by hand}\label{app:cone}

The framework assigns the Lorentz signature to the $\alpha=2$ closure (\S\ref{sec:galilean}). Nothing in this appendix describes the $\alpha\to1^-$ edge. It records what follows if a light cone is nevertheless imposed on the edge, as one ingredient of an eventual crossover to $\alpha=2$.

\emph{Setting.} Let the spatial profile depend on the lag and be supported in $\lvert y\rvert\le c\tau$. Take it self-similar, so that it is a superposition of thin shells of speed $u\in[0,c]$ with weight $w(u)$, and set $v_{\rm eff}^2=\int u^2w(u)\,du$. Hold $E_0\equiv\eps E_cr_1$ fixed. On a plane wave the field equation reads
\[
\alpha=\eps r_1e^{-i\sigma\alpha\pi/2}\tau_c^{-\alpha}m_\alpha(E,k),
\qquad
m_\alpha=\frac1{\Gamma(\alpha)}\int_0^\infty\tau^{\alpha-1}e^{iE\tau/\hbar}\hat g_\tau(k)\,d\tau .
\]
For a thin shell of speed $v$, with $s=-iE/\hbar$,
\[
m_\alpha=\frac{(s-ikv)^{1-\alpha}-(s+ikv)^{1-\alpha}}{2ikv(\alpha-1)} .
\]

\emph{Results, under the imposed cone.}
\begin{enumerate}
\item On $E>\hbar kc$ the symbol's phase cancels for $\sigma=+1$, so the cone forces the opposite orientation to the Galilean edge. The measured residual is $\lvert{\rm Im}\rvert/\lvert\cdot\rvert\le7\times10^{-15}$.
\item As $\alpha\to1^-$, $E(k)=\hbar kv\coth(\hbar kv/E_0)$, with $m=3E_0/(2v^2)$ and $0<v_g<v$. For a general self-similar profile, $m=3E_0/(2v_{\rm eff}^2)$.
\item Consequently $E_0\le\tfrac23mc^2$.
\item For $\alpha<1$: $E(0)=E_c(E_0/\alpha E_c)^{1/\alpha}$, $m_\alpha=3E(0)/((1+\alpha)v^2)$, and a cutoff $\hbar kv\approx\tfrac12E_c\big(E_0/(\alpha\eps E_c)\big)^{1/\alpha}$ that recedes as $\eps\to0$.
\end{enumerate}
The derivations are those of the plane-wave analysis above, applied to $m_\alpha$. They are kept here, with the numbers below, only as the record of an $\alpha=2$ ingredient.

\begin{table}[h]
\centering\small
\begin{tabular}{cccccc}
\toprule
$k$ & root & $kv\coth(kv/E_0)$ & $E_0+k^2/2m$ & rel.\ kinetic dev. & $v_g/v$\\
\midrule
$0.01$ & $1.000033$ & $1.000033$ & $1.000033$ & $-6.7\times10^{-6}$ & $0.007$\\
$0.10$ & $1.003331$ & $1.003331$ & $1.003333$ & $-6.7\times10^{-4}$ & $0.067$\\
$1.00$ & $1.313035$ & $1.313035$ & $1.333333$ & $-6.1\times10^{-2}$ & $0.589$\\
$3.00$ & $3.014909$ & $3.014909$ & $4.000000$ & $-0.33$ & $0.975$\\
\bottomrule
\end{tabular}
\caption{Imposed cone, thin shell, $v=E_0=1$. Not the $\alpha\to1^-$ edge.}
\end{table}

\begin{thebibliography}{99}
\bibitem{CAL} S.~Leizerman, \emph{The Causal Accumulation Law: The Four Corners of $\alpha$, Algebraic Closure, and the Standard Model}, LIOR Research Initiative (2026).
\bibitem{Skeleton} S.~Leizerman, \emph{The $\alpha=1$ Unitary Skeleton and the Emergence of Schr\"odinger Evolution in the $\delta$-Kernel Limit of the Causal Accumulation Law}, LIOR Research Initiative (2026).
\bibitem{Born} S.~Leizerman, \emph{An Axiomatic Derivation of the Born Rule at the $\delta$-Kernel Limit and Closure Order $\alpha=2$ of the Causal Accumulation Law}, LIOR Research Initiative (2026).
\bibitem{SKM} S.~G.~Samko, A.~A.~Kilbas, O.~I.~Marichev, \emph{Fractional Integrals and Derivatives: Theory and Applications}, Gordon and Breach (1993).
\bibitem{Podlubny} I.~Podlubny, \emph{Fractional Differential Equations}, Academic Press (1999).
\bibitem{GKMR} R.~Gorenflo, A.~A.~Kilbas, F.~Mainardi, S.~V.~Rogosin, \emph{Mittag-Leffler Functions, Related Topics and Applications}, Springer (2014).
\bibitem{Lubich} C.~Lubich, Convolution quadrature and discretized operational calculus~I, Numer.\ Math.\ \textbf{52}, 129--145 (1988).
\bibitem{Stone} M.~H.~Stone, On one-parameter unitary groups in Hilbert space, Ann.\ of Math.\ \textbf{33}, 643--648 (1932).
\bibitem{RS1} M.~Reed, B.~Simon, \emph{Methods of Modern Mathematical Physics I: Functional Analysis}, Academic Press (1980).
\bibitem{Kato} T.~Kato, \emph{Perturbation Theory for Linear Operators}, Springer (1995).
\end{thebibliography}

\end{document}
